\section{Introduction}\label{sec:introduction}
A cell-bounded contingency table is a nonnegative integer matrix whose
row sums, column sums, and individual entry bounds are prescribed.
For positive integers $m,n$, margins
$r\in\mathbb Z_{\ge0}^m$ and $c\in\mathbb Z_{\ge0}^n$ with equal total,
and bounds $b\in\mathbb Z_{\ge0}^{m\times n}$, write
\begin{equation}\label{eq:feasible-set}
\begin{aligned}
 \Omega(r,c,b)&=\left\{X\in\mathbb Z_{\ge0}^{m\times n}:
 \sum_{j=1}^n X_{ij}=r_i,\quad
 \sum_{i=1}^m X_{ij}=c_j,\quad X_{ij}\le b_{ij}\right\},\\
 Z(r,c,b)&=|\Omega(r,c,b)|.
\end{aligned}
\end{equation}
Each table is counted once. In particular, $b_{ij}=0$ forbids an entry;
no positivity condition is imposed on the support or margins.
Our concern is the binary encoding model, in which a capacity can be
exponentially larger than its description. Both $m$ and $n$ are part of
the input.

\begin{theorem}\label{thm:main}
There is one randomized algorithm with the following guarantee. Given
$(r,c,b)$ as above and rational $\varepsilon,\delta\in(0,1)$, it outputs
a nonnegative rational $\widehat Z$ such that
\[
 \P\bigl((1-\varepsilon)Z(r,c,b)\le \widehat Z
                  \le(1+\varepsilon)Z(r,c,b)\bigr)\ge1-\delta.
\]
If $\Omega(r,c,b)$ is empty, the output is zero on every execution.
The algorithm uses independent unbiased bits, and one fixed polynomial
in $L$, $\varepsilon^{-1}$, and $\log(\delta^{-1})$ bounds its total bit
operations on every execution, where $L$ is the total binary input
length, including the accuracy parameters. No counting or sampling
oracle is required.
\end{theorem}

The theorem allows arbitrary cell bounds and margins without a fixed
number of rows or columns, a lower bound on positive capacities, or a
density or balance assumption. Its running-time guarantee includes
preprocessing, random-bit generation, rational arithmetic, and output,
even on executions in the failure event. The polynomial bounds in the
proof are deliberately generous; the result is a uniform complexity
guarantee rather than a practical running-time estimate.

\subsection{History and significance}
The enumeration of fixed-margin arrays has long been studied in
statistics and combinatorics; Diaconis and Gangolli~\cite{DG1995}
survey its applications and exact and approximate methods.
Ordinary contingency tables are a special case of
\eqref{eq:feasible-set}, since the margins imply
$X_{ij}\le\min\{r_i,c_j\}$. Even their exact enumeration is difficult:
Dyer, Kannan, and Mount~\cite[Theorem~1]{DKM1997} proved
$\#\mathrm P$-completeness for tables with only two rows.
Approximate counting can nevertheless be efficient under much broader
conditions than exact counting.

For ordinary tables, Dyer, Kannan, and Mount~\cite{DKM1997}
developed geometric sampling and counting methods for sufficiently
large margins. Morris~\cite{Morris2002} weakened the required margin
bounds by refining the expanded-polytope and rounding argument.
Dyer and Greenhill~\cite{DG2000} gave polynomial-time approximate
counting and sampling for arbitrary two-row margins. Cryan and
Dyer~\cite{CD2003} obtained an FPRAS for any fixed number of rows by
combining dynamic programming and volume estimation; Dyer~\cite{Dyer2003}
subsequently improved the fixed-row algorithm using dynamic programming.
Gopalan, Klivans, Meka, \v{S}tefankovi\v{c}, Vempala, and
Vigoda~\cite[Theorem~1.3]{GKM2011} later gave a deterministic FPTAS
for a fixed number of rows.
These results already permit binary-encoded margins. Their restrictions
concern the number of rows or the sizes of the margins, rather than a
unary encoding convention.

Individual caps connect the problem to several other counting models.
When all caps are zero or one, a table is a subgraph of the allowed
bipartite graph with prescribed vertex degrees. Jerrum, Sinclair, and
Vigoda~\cite[Corollary~8.1]{JSV2004} give an FPRAS for this entire
special case as a consequence of their permanent algorithm; all-one
margins give perfect matchings. Bez\'akov\'a, Bhatnagar, and
Vigoda~\cite{BBV2007} developed a direct annealing algorithm for
binary tables, also allowing arbitrary forbidden cells.
More generally, capped tables are
integer flows in a bipartite network. Cryan, Dyer, and
Randall~\cite{CDR2010} gave an FPRAS for integral flows with sufficiently
large tight capacities, and for cell-bounded tables with a fixed number
of rows and arbitrary positive margins and caps. They identified
unrestricted approximation as an open problem in their 2010 paper.
Their tight-capacity reduction and maximum-capacity spanning-tree
representation are direct antecedents of our large-cell construction.
Here a forest and soft constraints accommodate arbitrary mixtures of
small and large bounds, with both dimensions variable and zero bounds
permitted.

Convex optimization gives a complementary family of estimates.
Barvinok~\cite[Theorem~1.3]{Barvinok2009} bounds ordinary weighted
table counts, including forbidden cells, within a factor
$N^{O(m+n)}$, where $N$ is the total margin. Barvinok, Luria,
Samorodnitsky, and Yong~\cite{BLSY2010} obtain randomized relative
approximations with complexity $N^{O(\log N)}$ for specified
smooth-margin classes. This scale is quasipolynomial in the numerical
total, rather than in its binary length. Br\"and\'en, Leake, and
Pak~\cite[Theorem~2.1]{BLP2021} give capacity bounds for arbitrary
individual caps, with explicit multiplicative factors depending on
the margins and dimensions. Their bounds extend earlier estimates
using Lorentzian polynomials. A subsequent preprint of September 30,
2026, by Leake and Mohammadi Yekta~\cite[Theorems~1.1 and~1.3]{LMY2026}
extends capacity-based estimates to finite nonnegative integer fibers
of totally unimodular systems, retaining an explicit multiplicative
factor depending on the system's dimensions and right-hand side.
Theorem~\ref{thm:main} supplies
any requested relative tolerance in polynomial binary time, without
a smoothness assumption.

The geometric problem here is to count all integer points of a
transportation polytope. These need not all be vertices: the
$2\times2$ tables with every margin equal to two form a segment
with three integer points and two vertices. Guo and
Jerrum~\cite[Proposition~4.2]{GJ2023} give an FPRAS for the vertices
of promised $0/1$ polytopes $\{x:0\le Ax\le1\}$ defined by network
matrices; this unit-bound vertex result has a different scope.

Counting also provides approximate sampling through the
self-reduction principle of Jerrum, Valiant, and
Vazirani~\cite[Theorem~6.3]{JVV1986}. Section~\ref{sec:flows}
implements this principle by bisecting cell intervals, with exact
feasibility tests at every branch. For an explicitly given directed
multigraph with finite integer arc bounds and prescribed vertex
balances, it yields both an FPRAS for integer arc-value vectors and
a sampler that always returns a feasible vector when one exists.
The latter has total-variation error at most $\tau$ and an
all-execution cost polynomial in the input length and $\tau^{-1}$.
A separate companion~\cite[Theorem~1.1]{OpenAIExactTables2026}
proves exact uniform sampling in polynomial expected bit time for
ordinary tables with arbitrary margins and no input cell bounds or
forbidden positions. Its bounded approximate sampler has cost
polynomial in the input length and the number of accuracy bits. Those stronger sampling
guarantees concern the uncapped setting; the flow consequence here
has the total-variation and cost bounds just stated.


\subsection{From tables to weighted transversals}
The proof separates numerical size from combinatorial complexity.
Large binary bounds cannot be expanded into that many binary choices.
We instead sum over large coordinates through a weighted completion
problem, and expand only coordinates with polynomially bounded ranges.
The resulting finite-state argument must preserve one unit of mass per
table despite the multiplicities introduced by this expansion.

\paragraph{Large entries and two views of a small entry.}
Flow feasibility first tightens each cell to its attainable interval.
Subtract its minimum, so that the new endpoints are $0$ and $b_s$;
each endpoint is attained by some feasible table. Remove zero-width
cells and separate the remaining cells at a polynomial threshold $W$.
In the bipartite graph of large cells, choose a maximum-capacity
spanning forest and one root in each component. The large cells outside
the forest are free coordinates. Once their values and the contributions
of small cells to each incident margin have been specified, leaf
elimination uniquely determines the tree entries from the nonroot
margin equations.

For a small cell $s$, temporarily allow two counts $x_s,y_s$ between
$0$ and $b_s$: its row contribution is $x_s$, and its column
contribution is $b_s-y_s$. The views agree when $x_s+y_s=b_s$.
Writing $p=\sum_{s\text{ small}}b_s$, we also consider unbalanced
profiles satisfying only $\sum_s(x_s+y_s)=p$. Extending the weights
to these profiles supplies the auxiliary states and algebraic tests
used below. For every such
profile, sum over the free large entries in their integer intervals.
Penalize the resulting deviations from root margins and violations of
tree-cell bounds by a nonnegative quantity $D$. This penalty vanishes
exactly when those remaining constraints hold. A maximum-capacity
forest bounds its slopes in normalized free coordinates independently
of the numerical large capacities.

For a polynomially bounded integer $H$, let $h_j(x,y)$ be the sum of
$2^{-jD/H}$ over these free assignments, for $0\le j\le H$, and
let $C(j)$ be the sum of $h_j$ over balanced small profiles. Thus
$C(0)$ is an explicit product of interval sizes. Every feasible table
contributes one to $C(H)$, while invalid completions contribute
additional positive mass. Log-concavity of each tree-coordinate marginal,
together with its attainable endpoints, controls violations close to
a tree boundary; exponential decay controls more distant violations.
This proves $Z\le C(H)\le(1+\epsilon/32)Z$. The use of scaled free
coordinates and soft constraints builds on the geometric table methods
of Dyer, Kannan, and Mount and Morris~\cite{DKM1997,Morris2002}.

\paragraph{Binary encodings and two extensions of their weights.}
Replace a small cell of bound $b_s$ by $b_s$ labeled pairs, each
containing a row element and a column element. There are $p$ pairs
altogether. A \emph{transversal} selects one element from each pair.
Its selected row and column counts give a balanced profile $(x,y)$.
The number of transversals with that profile is
$\prod_s\binom{b_s}{x_s}$, so assigning each the weight
$h_j(x,y)/\prod_s\binom{b_s}{x_s}$ gives mass $h_j(x,y)$ to that
profile and total mass $C(j)$ to all transversals.

The transport argument uses an enlarged chain that also permits $p$-element sets
with one empty pair and one full pair. Such a state has \emph{defect
type} $(i,l)$ when pair $i$ is empty and pair $l$ is full.
Two factorial formulas extend the transversal weights to positive
weights $f_j$ and $\widetilde f_j$ on all $p$-element sets. They
agree on transversals and are within a polynomial factor of each other
on the permitted defect types. Their algebraic roles differ. For $f_j$,
matrices of weights obtained by omitting two elements from a fixed
$(p+2)$-element set have at most one positive eigenvalue. For
$\widetilde f_j$, the corresponding statement holds when adding two
elements to a fixed $(p-2)$-element set. These quadratic tests
belong to the Lorentzian framework of Br\"and\'en and
Huh~\cite{BH2020}; we prove the specific integer summation and binary
contraction statements needed here. The omission property controls
transport between conditional transversal laws. The addition property
compares defect masses. Passing the latter inequalities between the
two weights incurs only polynomial losses.

\paragraph{The chain, its observables, and bounded computation.}
We simulate exchanges of one selected and one unselected element,
using $f_j$ and a multiplier for each defect type. As in the permanent
algorithm of Jerrum, Sinclair, and Vigoda~\cite[Section~3]{JSV2004},
the multipliers balance the type masses and are learned from observed
type frequencies during annealing. The paired transport, variance,
and trace analysis adapts the common-bases construction
of~\cite[Sections~4--6]{OpenAICommonBases2026}, with the two distinct
weights supplying its two signature inputs.

The quantities to estimate are the adjacent ratios $C(j+1)/C(j)$
and the defect-type probabilities. A variance bound controls these
observables even though it does not bound arbitrary fluctuations
within each defect type. Restricted variance estimates also underlie
the permanent algorithm of Chen, Vigoda, and Yang~\cite{CVY2026}.
Here we additionally control the chain's \emph{transversal trace},
which records successive returns to the transversal class. Its mixing
and return-cost bounds provide fresh starts for estimating the ratios;
their product recovers $C(H)$ from $C(0)$.

Finally, the completion weights $h_j$ are evaluated approximately by
a walk on polynomially many bins per free coordinate, with a correction
for the different numbers of integers in the bins. The product grid
itself can be exponentially large; only polynomially many transitions
are simulated. Fresh weight estimates and finite-bit choices are
coupled to the exact-weight analysis at each adaptive request. Caps on
trace restarts and all numerical subroutines give a polynomial work
bound on every execution, while the accuracy analysis bounds the
probability of an inaccurate estimate or a capped restart.

\subsection{Organization and conventions}
Section~\ref{sec:scales} constructs the forest coordinates and softened
counts. Section~\ref{sec:integer-signatures} proves the integer summation
properties and controls contamination by invalid completions.
Section~\ref{sec:signatures} introduces the binary weights and proves
their signatures and defect inequalities. Section~\ref{sec:transport}
derives transport, observable variance, and trace estimates.
Section~\ref{sec:evaluation} implements approximate evaluation of the
completion weights. Section~\ref{sec:algorithm} specifies
annealing, implements it using fair bits, and proves
Theorem~\ref{thm:main}. Section~\ref{sec:flows} gives the counting and
total-variation sampling consequences for bounded integral network flows.

All logarithms with a subscript $2$ are binary; $\log$ in exponential
estimates is natural. Sums over edges in Dirichlet forms are over
unordered edges unless stated otherwise. A symmetric matrix has
\emph{at most one positive eigenvalue} when its positive inertia index
is at most one, with multiplicity. This property concerns quadratic
forms and is preserved by restriction and limits; preservation under
pair summation will be proved, not assumed. The large matrices,
conditioning trees, and transport flows used in these proofs are not
constructed by the algorithm.
